\documentclass[11pt]{article}
\usepackage[margin=1in]{geometry}
\usepackage[T1]{fontenc}
\usepackage{lmodern,amsmath,amssymb,amsthm,mathtools,hyperref}
\setlength{\emergencystretch}{3em}
\newtheorem{theorem}{Theorem}[section]
\newtheorem{lemma}[theorem]{Lemma}
\newtheorem{corollary}[theorem]{Corollary}
\newtheorem{proposition}[theorem]{Proposition}
\newcommand{\R}{\mathbb R}
\newcommand{\tr}{\operatorname{tr}}
\newcommand{\Sym}{\operatorname{Sym}}
\newcommand{\Skew}{\operatorname{Skew}}
\title{Tensor powers destabilize nonflat critical points\\of entrywise orthogonal $p$-energies}
\author{Joshua Mark\thanks{ORCID: \href{https://orcid.org/0009-0003-5905-2382}{0009-0003-5905-2382}. AI tools assisted with research, derivation, adversarial review, computation and exposition.}}
\date{20 November 2026 --- research draft; priority not established}
\begin{document}
\maketitle
\begin{abstract}
For every finite real $p>2$, we prove that a nonflat critical point of
the entrywise $p$-energy on the real orthogonal group acquires a negative
tangent-Hessian direction at some finite self-tensor power. The proof
uses symmetric and skew-symmetric trace averages and strict Chebyshev
inequality. It gives an explicit bound on the number of tensor factors
and a factored descent direction. Consequently the orthogonal matrices
whose every tensor power is a local minimum are exactly the scaled
Hadamard matrices. This does not prove that such matrices exist in
every admissible dimension: tensor lifting changes the dimension, and
no descent back to the original dimension is established. Novelty is
not asserted pending a comprehensive literature review.
\end{abstract}

\section{Main statement}
For $n\ge2$ and $2<p<\infty$, define
\[
 F_p(U)=\sum_{i,j=1}^n |u_{ij}|^p\qquad(U\in O(n)).
\]
Call $U$ flat if all its entries have magnitude $1/\sqrt n$.
Flatness is equivalent to $\sqrt nU$ being a Hadamard matrix.
Jensen's inequality and $\sum u_{ij}^2=n$ give
\[
 F_p(U)\ge n^{2-p/2},
\]
with equality exactly when $U$ is flat. Entrywise norm optimization
and its differential framework have been studied in the almost-Hadamard
literature~\cite{BN}. Our assertion concerns stability under arbitrary
\emph{self-tensor powers}, not just the landscape at one fixed order.

\begin{theorem}[Finite tensor destabilization]\label{thm:main}
Let $U\in O(n)$ be a nonflat critical point of $F_p$. Put
\[
 q=\sum_{ij}|u_{ij}|^p,\qquad
 r=\sum_{ij}|u_{ij}|^{p-2},\qquad
 b=\frac{r+q}{(n+1)q}.
\]
Then $0<b<1$. For some positive integer $k$, the critical point
$U^{\otimes k}\in O(n^k)$ has a nonzero skew-symmetric tangent
generator $X$ such that
\[
 \left.\frac{d^2}{dt^2}F_p(U^{\otimes k}e^{tX})\right|_{t=0}<0.
\]
One may take
\begin{equation}\label{eq:kbound}
 k=2+\left\lceil\frac{\log(p-1)}{\log(1/b)}\right\rceil.
\end{equation}
The direction can be chosen in the factored form
$X=Y\otimes Z^{\otimes(k-1)}$, with $Y$ skew-symmetric and $Z$
symmetric. Both may be chosen from elementary matrix bases.
\end{theorem}

The result allows zero entries. The function $|x|^p$ is $C^2$ on the
whole real line for every $p>2$, with second derivative zero at zero.
No spectral positivity assumption on the critical-point matrix below
is needed.

\section{Derivative and trace identities}
Write
\[
 W=\bigl(|u_{ij}|^{p-2}u_{ij}\bigr)_{ij},\qquad M=U^{\mathsf T}W.
\]
Differentiation in skew directions shows that criticality is equivalent
to $M$ being symmetric. Its trace is $q$.
For any real $n\times n$ matrix $V$, set
\[
 T(V)=\sum_{ij}|u_{ij}|^{p-2}(UV)_{ij}^2,\qquad
 S(V)=\tr(M V^{\mathsf T}V).
\]
We always have $T(V)\ge0$, but individual values of $S$ may be negative.
For a skew-symmetric $V$,
\begin{equation}\label{eq:hessian}
 \left.\frac{d^2}{dt^2}F_p(Ue^{tV})\right|_0
 =p\bigl((p-1)T(V)-S(V)\bigr).
\end{equation}
Indeed the unreduced second derivative is
$p(p-1)\sum |u_{ij}|^{p-2}(UV)_{ij}^2+
p\tr(W^{\mathsf T}UV^2)$; symmetry of $M$ and
$V^2=-V^{\mathsf T}V$ give the formula.

\begin{lemma}[Symmetric and skew trace sums]\label{lem:traces}
Let $\mathcal B_+$ and $\mathcal B_-$ be Frobenius-orthonormal bases
of the symmetric and skew-symmetric matrices respectively. Then
\begin{align*}
 \sum_{Z\in\mathcal B_+}T(Z)&=(r+q)/2,&
 \sum_{Z\in\mathcal B_+}S(Z)&=(n+1)q/2,\\
 \sum_{Y\in\mathcal B_-}T(Y)&=(r-q)/2,&
 \sum_{Y\in\mathcal B_-}S(Y)&=(n-1)q/2.
\end{align*}
\end{lemma}
\begin{proof}
For the unit row vector $u_i$ of $U$, the squared norms of the
symmetric and skew projections of $u_i^{\mathsf T}e_j^{\mathsf T}$
are $(1+u_{ij}^2)/2$ and $(1-u_{ij}^2)/2$. Equivalently, direct
summation over $E_{ii}$ and $(E_{ij}\pm E_{ji})/\sqrt2$ gives
\[
 \sum_{Z\in\mathcal B_\pm}(UZ)_{ij}^2
 =\frac{1\pm u_{ij}^2}{2}.
\]
Multiply by $|u_{ij}|^{p-2}$ and sum to obtain the $T$ identities.
The same elementary bases satisfy
\[
 \sum_{Z\in\mathcal B_\pm}Z^{\mathsf T}Z
 =\frac{n\pm1}{2}I_n.
\]
Taking their trace product with $M$ proves the $S$ identities.
\end{proof}

\begin{lemma}[Strict moment inequality]\label{lem:cheb}
For a nonflat $U$, $q\le r<nq$.
\end{lemma}
\begin{proof}
The first inequality follows from $|u_{ij}|\le1$. For the second,
list the $n^2$ squared entries as $x_1,\ldots,x_{n^2}$ and put
$a=(p-2)/2>0$. Then $\sum x_i=n$, $\sum x_i^a=r$, and
$\sum x_i^{a+1}=q$. Strict Chebyshev inequality follows directly from
\[
 n^2q-nr
 =\sum_{i<j}(x_i-x_j)(x_i^a-x_j^a)>0.
\]
Every summand is nonnegative; at least one is positive because the
$x_i$ are not all equal. This argument includes zero entries.
\end{proof}

\begin{lemma}[Ratio selection without positivity of every denominator]
Suppose finitely many pairs $(T_i,S_i)$ have $T_i\ge0$ and
$\sum S_i>0$. Some $i$ has $S_i>0$ and
\[
 T_i/S_i\le\frac{\sum T_j}{\sum S_j}.
\]
\end{lemma}
\begin{proof}
Set the displayed aggregate ratio to $c\ge0$. If all positive-$S_i$
terms had $T_i>cS_i$, while nonpositive-$S_i$ terms satisfy
$T_i\ge0\ge cS_i$, summing would contradict the definition of $c$.
\end{proof}

\section{Proof of tensor destabilization}
Apply the trace sums and ratio-selection lemma to choose $Y\in
\mathcal B_-$ and $Z\in\mathcal B_+$ with positive $S$-values and
\[
 \alpha:=\frac{T(Y)}{S(Y)}
 \le\frac{r-q}{(n-1)q}<1,\qquad
 \beta:=\frac{T(Z)}{S(Z)}\le b<1.
\]
The upper bound on $\alpha$ is nonnegative, and the choices exist
because $n\ge2$ and $q>0$.

At the tensor power, $W_k=W^{\otimes k}$ and $M_k=M^{\otimes k}$.
Hence $U^{\otimes k}$ is still critical. For
$X=Y\otimes Z^{\otimes(k-1)}$, which is nonzero and skew-symmetric,
tensor factorization of the sums and traces gives
\[
 T_k(X)=T(Y)T(Z)^{k-1},\qquad
 S_k(X)=S(Y)S(Z)^{k-1}.
\]
Its Hessian is therefore
\begin{equation}\label{eq:tensorhessian}
 pS(Y)S(Z)^{k-1}
 \bigl((p-1)\alpha\beta^{k-1}-1\bigr).
\end{equation}
If $\alpha=0$, $k=1$ already works. Otherwise, choose $k$ as in
\eqref{eq:kbound}. Since $\alpha<1$, $\beta\le b<1$, and
$k-1>\log(p-1)/\log(1/b)$, the bracket is negative. This proves
Theorem~\ref{thm:main}.

\begin{corollary}[Characterization by tensor stability]
For $p>2$, the following are equivalent for $U\in O(n)$:
\begin{enumerate}
\item Every positive self-tensor power of $U$ is a local minimum of
the entrywise $p$-energy in its ambient orthogonal group.
\item Every positive self-tensor power is critical with positive
semidefinite tangent Hessian.
\item $U$ is flat, so $\sqrt nU$ is a Hadamard matrix.
\end{enumerate}
\end{corollary}
\begin{proof}
The first implication is the second-order necessary condition. The
theorem proves the second implication to flatness. If $U$ is flat,
every tensor power is flat and attains the Jensen lower bound globally.
\end{proof}

\begin{corollary}[Strict eventual submultiplicativity]
Let $m_p(n)=\min_{U\in O(n)}F_p(U)$. If there is no Hadamard matrix
of order $n$, then for some positive $k$,
\[
 m_p(n^k)<m_p(n)^k.
\]
Equivalently, equality for every $k$ occurs precisely when order $n$
admits a Hadamard matrix.
\end{corollary}
\begin{proof}
Compactness gives a minimizer $U$, which is critical and nonflat under
the hypothesis. The negative Hessian direction strictly decreases
$F_p$ near $U^{\otimes k}$, whose value is $m_p(n)^k$. Conversely a
flat matrix and all its tensor powers attain the Jensen lower bounds.
\end{proof}

\begin{corollary}[Rectangular row-orthonormal frames]
The tensor-destabilization theorem and tensor-stability characterization
also hold on
\[
 \operatorname{St}(d,n)=\{U\in\R^{d\times n}:UU^{\mathsf T}=I_d\},
 \qquad 1\le d\le n,\quad n\ge2.
\]
Here flat means every entry has magnitude $1/\sqrt n$, so the scaled
matrix is a rectangular partial Hadamard matrix. The tensor powers have
shape $d^k\times n^k$, and the same bound~\eqref{eq:kbound} uses the
number $n$ of columns.
\end{corollary}
\begin{proof}
Right multiplication by $O(n)$ acts transitively on this Stiefel
manifold, so criticality is again equivalent to symmetry of the
$n\times n$ matrix $M=U^{\mathsf T}W$. All rows still have norm one,
so both trace-sum identities are unchanged. For strict Chebyshev, there
are now $dn$ squared entries summing to $d$; the same pair sum is
$dnq-dr>0$ unless every entry is flat. Thus $r<nq$ still holds and
the ratio and tensor arguments apply without change. Positive
$S(Y),S(Z)$ ensure that $UY$ and $UZ$ are nonzero. Indeed, if
$UV=0$ for a symmetric or skew $V$, then $VU^{\mathsf T}=0$ and
$V$ annihilates the range of $M$, implying $S(V)=0$.
The factored actual tangent is therefore nonzero. Finally flat row frames minimize the energy
globally, by Jensen separately on each unit row, with value
$d n^{1-p/2}$.
\end{proof}

\section{A local minimum that becomes unstable after two copies}
For $p=4$, let $A=2J_5/5-I_5$. The earlier note proves that
$A\otimes(H_k/\sqrt k)$ is a strict local minimum with value $29/25$
for any Hadamard factor. At $A$,
\[
 M=\frac7{25}I_5-\frac6{125}J_5.
\]
Let $Y=E_{01}+E_{12}-E_{02}$ be a skew triangle circulation. Let $Z$
be symmetric, with $Z_{01}=1$, $Z_{12}=-1$, $Z_{23}=1$,
$Z_{30}=-1$, the transposed entries equal, and all other entries zero.
Both have row sums zero and zero diagonal, and
$\|Y\|_F^2=6$, $\|Z\|_F^2=8$. Thus
\[
 T(Y)=\tfrac4{25}\|Y\|_F^2,\quad
 S(Y)=\tfrac7{25}\|Y\|_F^2,
\]
with the same identities for $Z$. Consequently at $A\otimes A$ the
direction $Y\otimes Z$ has quartic curvature
\[
 4\cdot\frac{49}{625}\cdot48
 \left(3\frac{16}{49}-1\right)=-\frac{192}{625}.
\]
Tensoring with a flat factor preserves this displayed curvature when
the tangent generator is tensored with an identity. In particular a
strict local minimum at order $20$ becomes unstable after taking its
tensor square at order $400$. This is consistent with stability under
flat tensor factors: self-tensoring inserts a second nonflat factor.

\section{Scope, checks and priority}
The proof is uniform in all real finite $p>2$, not a claim inferred
from numerical Hessian spectra. The accompanying standard-library
Python checker verifies the trace formulas, stationarity, ratio choices,
and factored negative curvatures with exact rational arithmetic for
integer exponents $p=3,4,5,6,8$ and several critical matrices. It also
checks the two-copy curvature directly at order $25$.

A separate mathematical review checked the trace formulas, strictness
with zeros, possibly indefinite $M$, and the tensor factorization.
This is not a Lean proof or a substitute for external peer review.
Targeted primary-literature searches on almost-Hadamard matrices,
tensor stability, local norm extrema and tensor-power descent did not
locate this precise statement. That does not establish priority.

Most importantly, this result does not show that a tensor-stable local
minimum exists in any prescribed dimension. Tensor descent occurs in
$O(n^k)$, and no mechanism here transfers it back to $O(n)$. Therefore
the characterization is an optimization result and a research lead,
not a proof of the Hadamard conjecture.

\begin{thebibliography}{1}
\bibitem{BN} T. Banica and I. Nechita,
\emph{Almost Hadamard matrices: the case of arbitrary exponents},
Discrete Appl. Math. 161 (2013), 2367--2379.
\url{https://arxiv.org/abs/1211.2669}.
\end{thebibliography}
\end{document}
